\uuid{aBOQ}
\exo7id{4098}
\titre{exo7 4098}
\auteur{quercia}
\organisation{exo7}
\datecreate{2010-03-11}
\isIndication{false}
\isCorrection{true}
\chapitre{Equation différentielle}
\sousChapitre{Equations différentielles linéaires}
\cpge{mpsi-mp2i.primitives-equations-differentielles.c, pcsi.techniques-calcul-primitives.c, ptsi.techniques-calcul-primitives.c, mpsi-mp2i.limites-continuite.c, pcsi.limites-continuite.c, ptsi.limites-continuite.c, bcpst1.equations-differentielles-lineaires.a, mp-mpi.equations-differentielles-lineaires.b, psi.equations-differentielles, pt.equations-differentielles.a}
\module{Analyse}
\niveau{L2}
\difficulte{}

\contenu{
\question{
Soit $\lambda > 0$ et $y$ une solution de
$y'' + \left(1+\frac\lambda{t^2}\right)y = 0$.
Montrer que pour tout $a \in \R$, $y$ a un zéro dans l'intervalle $]a,a+\pi[$.
(\'Etudier $z = y'\varphi - y\varphi'$ où $\varphi(t) = \sin(t-a)$)
}
\reponse{
$z' = -\frac\lambda{t^2}\sin(t-a)y(t)$ donc si $y$ ne s'annule pas
sur $]a,a+\pi[$, alors $z$ est strictement monotone sur $[a,a+\pi]$.
Mais $z(a+\pi)-z(a) = y(a+\pi)+y(a)  \Rightarrow $ contradiction de signe.
}
}
